% !TeX root = ../../Book2.tex
\section{应用}
\label{b2:sec:2203}

应用乘子定理时，先判断符号阶数，再选取匹配的函数空间。零阶符号可保持 \(L^p\)，负阶符号在不同频率区间可能改善正则性或改变可积指数，正阶符号则需输入具有相应导数。

\subsection{Riesz 变换}
\label{b2:subsec:220301}

Riesz 变换的符号 \(-\mathrm i\xi_j/|\xi|\) 满足任意阶的尺度导数条件，故\chapref{b2:ch:18}的 \(L^p\) 定理也可由 Mihlin 定理再次得到。两条路线揭示不同方面：奇异积分证明直接看到零平均和核尾部；乘子证明利用频率光滑性，以 \(L^2\) 自动建立算子起点。

\ProseParagraphBreak
还可处理任意光滑零次齐次符号 \(m\)，因为单位球面上的有限阶导数有界。零次齐次性并不要求球面平均为零；常数符号的核有一个原点质量。只有在提取原点外的奇异核时，才需要辨认其球面抵消，不能把抵消条件错误地加到所有零阶符号上。

\subsection{分数幂}
\label{b2:subsec:220302}

记
\[
 J^sf=\mathcal F^{-1}\bigl[(1+4\pi^2|\xi|^2)^{s/2}\widehat f\bigr],
 \qquad (-\Delta)^{s/2}f
 =\mathcal F^{-1}\bigl[(2\pi|\xi|)^s\widehat f\bigr].
\]
前者是不齐次分数阶算子，后者是齐次分数阶算子。对 \(s>0\)，\(J^{-s}\) 的符号满足 Mihlin 条件，故在 \(L^p\) 上有界；\((-\Delta)^{-s/2}\) 在零频发散，应在适当齐次空间或以\chapref{b2:ch:19}的分数积分解释，不能据此声称它在同一个 \(L^p\) 上有界。

\ProseParagraphBreak
对 \(1<p<\infty\)，定义 Bessel 位势空间
\[
 H^{s,p}=\{f\in\mathcal S':J^sf\in L^p\},\qquad
 \|f\|_{H^{s,p}}=\|J^sf\|_p.
\]
这是完整空间，因为 \(J^{-s}\) 把 \(L^p\) 连续送入温和分布，\(L^p\) 极限可经它反解。连续性的核对可对任意测试函数 \(\psi\) 使用
\(|\langle J^{-s}g,\psi\rangle|\le\|g\|_p\|J^{-s}\psi\|_{p'}\)；
任意实 \(s\) 时 \(J^{-s}\psi\) 仍为 Schwartz 函数。

\begin{proposition}\label{b2:prop:22-bessel-square}
对非齐次分割 \(P_0,\Delta_j\)（\(j\ge1\)），有
\[
 \|J^sf\|_p\asymp
 \left\|\left(|P_0f|^2+\sum_{j\ge1}2^{2js}|\Delta_jf|^2\right)^{1/2}\right\|_p.
\]
因而\chapref{b2:ch:19}定义的 \(F^s_{p,2}\) 与 \(H^{s,p}\) 相同，范数等价。
\end{proposition}
\begin{proof}
设 \(g=J^sf\)。第 \(j\) 个加权块的符号为
\(2^{js}\varphi(2^{-j}\xi)(1+4\pi^2|\xi|^2)^{-s/2}\)。
在归一化环带上，它及有限阶导数一致有界，核具有\chapref{b2:ch:19}向量值平方函数证明使用的一致大小、光滑与 \(L^2(\ell^2)\) 界。于是该证明给出
\(\|(\sum2^{2js}|\Delta_jf|^2)^{1/2}\|_p\le C\|g\|_p\)。
低频块同样由一个 Schwartz 卷积核控制。

反向利用环带上等于 \(1\) 的伴随截断 \(\widetilde\varphi_j\) 及分解恒等式。
与测试函数 \(h\in\mathcal S\) 配对，把 \(J^sf\) 写成块的重构，每项中的 \(J^s\) 除以 \(2^{js}\) 后又是一个一致零阶环带符号。逐点 Cauchy–Schwarz 与\chapref{b2:ch:19}相同的向量值界得到
\[
 |\langle J^sf,h\rangle|
 \le C\left\|\left(|P_0f|^2+\sum2^{2js}|\Delta_jf|^2\right)^{1/2}\right\|_p
          \|h\|_{p'}.
\]
先对有限和证明，再由平方函数尾部的支配收敛及分布重构取极限；对偶性给出反向界。这里没有用习题中的未证明等价，也没有把一般 \(F^s_{p,q}\) 的理论混进来。
\end{proof}
当 \(s\) 为非负整数，导数符号除以 \(J^s\) 满足 Mihlin 条件。反向则以有限个局部符号分解或椭圆符号恒等式控制 \(J^s\)，得到通常 Sobolev 导数范数的等价。仅在 \(p=2\) 时，这种等价可以完全不经过乘子理论而直接由 Plancherel 得到。

% 扩写定位：b2:subsec:220302
整数阶的识别可以把所需符号明确写出。设 \(k\geq1\) 为整数，\(1<p<\infty\)。正向符号
\[
 \frac{(2\pi\mathrm i\xi)^\alpha}{(1+4\pi^2|\xi|^2)^{k/2}},
 \qquad |\alpha|\leq k,
\]
满足 Mihlin 条件，所以 \(\sum_{|\alpha|\leq k}\|\partial^\alpha f\|_p\leq C\|J^kf\|_p\)。
反向先在 \(|\xi|\leq2\) 取一个低频截断，再将 \(|\xi|\geq1\) 的球面分成有限个坐标锥：每片有某个 \(|\xi_j|\geq c|\xi|\)。以光滑齐次分割 \(\chi_j\) 写成
\[
 (1-\chi_0(\xi))(1+4\pi^2|\xi|^2)^{k/2}
 =\sum_{j=1}^d a_j(\xi)(2\pi\mathrm i\xi_j)^k,
\]
其中 \(a_j=(1-\chi_0)\chi_j(1+4\pi^2|\xi|^2)^{k/2}/(2\pi\mathrm i\xi_j)^k\)，在其支集附近定义并平滑延拓为零。分母在该锥上有统一下界，每个 \(a_j\) 满足零阶导数条件。低频符号 \(\chi_0(1+4\pi^2|\xi|^2)^{k/2}\) 的核为 Schwartz 函数。Mihlin 与 Young 给出
\[
 \|J^kf\|_p\leq C\left(\|f\|_p+\sum_{j=1}^d\|\partial_j^kf\|_p\right).
\]
初始等式在 Schwartz 类上成立，再通过分布与光滑逼近得到 \(H^{k,p}=W^{k,p}\)。这里连纯坐标的最高阶导数就已足够控制全部混合导数；它是乘子分解的结论，而不是可以在任意区域或任意边界条件下直接套用的代数事实。

\ProseParagraphBreak

分数阶算子还具有一个反映非局部性的实空间公式。对 \(0<s<2\)，设
\[
 A_{d,s}=\int_{\mathbb R^d}\frac{1-\cos(2\pi z_1)}{|z|^{d+s}}\,\mathrm dz,
 \qquad c_{d,s}=\frac{(2\pi)^s}{A_{d,s}}.
\]
原点处 \(1-\cos(2\pi z_1)\leq C|z|^2\)，无穷处它有界，故 \(0<A_{d,s}<\infty\)；这同时解释范围 \(0<s<2\)。
\begin{proposition}\label{b2:prop:22-fractional-kernel}
对 \(f\in\mathcal S\)，
\[
 (-\Delta)^{s/2}f(x)=\frac{c_{d,s}}2
 \int_{\mathbb R^d}\frac{2f(x)-f(x+y)-f(x-y)}{|y|^{d+s}}\,\mathrm dy.
\]
并且
\[
 \langle(-\Delta)^{s/2}f,f\rangle
 =\frac{c_{d,s}}2
 \iint\frac{|f(x)-f(y)|^2}{|x-y|^{d+s}}\,\mathrm dx\,\mathrm dy.
\]
\end{proposition}
\begin{proof}
对称二阶差在原点是 \(O(|y|^2)\)，远场由函数有界和核尾部可积控制，故第一式绝对收敛。先在 \(\epsilon<|y|<R\) 上作 Fourier 变换，得到乘子
\[
 c_{d,s}\int_{\epsilon<|y|<R}
 (1-\cos(2\pi y\cdot\xi))|y|^{-d-s}\,\mathrm dy.
\]
旋转及换元 \(z=|\xi|y\) 把全积分变成
\(c_{d,s}|\xi|^s A_{d,s}=(2\pi|\xi|)^s\)。可用非负积分的单调收敛，或用 Schwartz 配对的支配收敛取极限，证明第一式。第二式由交换 \(x,y\)、展开平方得到，等价地也可由 Plancherel 与同一个非负乘子积分推导；非负性允许 Tonelli。
\end{proof}
若 \(f\geq0\) 且紧支集，在支集外有
\((-\Delta)^{s/2}f(x)=-c_{d,s}\int f(y)|x-y|^{-d-s}\,\mathrm dy<0\)，除非 \(f=0\)。
因此这个导数一般不保持支集，也不保持非负性；它的二次型却非负。这三个性质并不矛盾：非负的是总能量，而不是每一点的算子输出。与整数 Laplace 算子的局部性相比，这个尾部正是分数阶模型的本质区别。

频率定义中的 Bessel 势也有一个不振荡的实空间核。设 \(s>0\)，定义
\[
 G_s(x)=\frac1{\Gamma(s/2)}
 \int_0^\infty t^{s/2-1}e^{-t}
 (4\pi t)^{-d/2}e^{-|x|^2/(4t)}\,\mathrm dt.
\]
这是把热核作正权平均，权的总质量为一。它与齐次 Riesz 势的区别恰在额外因子 \(e^{-t}\)：大时间被截去，因而不会保留齐次核的长距离幂尾。
\begin{proposition}\label{b2:prop:22-bessel-kernel}
有 \(G_s\geq0\)、\(\int G_s=1\)、\(\widehat G_s(\xi)=(1+4\pi^2|\xi|^2)^{-s/2}\)。故 \(J^{-s}f=G_s*f\)，并在全部 \(1\leq p\leq\infty\) 上满足 \(\|J^{-s}f\|_p\leq\|f\|_p\)。
当 \(0<s<d\) 时，\(G_s(x)\asymp |x|^{s-d}\) 于 \(0<|x|\leq1\)；当 \(s=d\) 时量级为 \(1+\log(1/|x|)\)；当 \(s>d\) 时在原点附近有界。无穷远处三种情形均有指数衰减。
\end{proposition}
\begin{proof}
非负性与 Tonelli 给出总质量一。热核 Fourier 变换为 \(e^{-4\pi^2t|\xi|^2}\)，对 \(t\) 的 gamma 积分得到所述符号，Young 即给出压缩界。

记 \(r=|x|\)。在 \(0<r\leq1\) 时，积分的 \(t\geq1\) 部分有统一有限上界。对 \(0<t<1\)，用 \(e^{-1}\leq e^{-t}\leq1\)，再换元 \(u=r^2/(4t)\)，得到与
\[
 r^{s-d}\int_{r^2/4}^\infty
 u^{(d-s)/2-1}e^{-u}\,\mathrm du
\]
可比的表达。若 \(s<d\)，积分趋于一个正的 gamma 常数；若 \(s=d\)，零点附近是 \(\int_{r^2/4}^1\mathrm du/u\)；若 \(s>d\)，它的下限奇性恰好抵消外面的 \(r^{s-d}\)，结果有界。下界可以分别在固定 \(u\) 区间或 \(r^2/4<u<1\) 中取出，故这些不仅是上界。

当 \(r\geq1\)，在指数中使用 \(t+r^2/(4t)\geq r\)，抽出 \(e^{-r/2}\)。剩下的指数 \(e^{-t/2-r^2/(8t)}\) 把 \(t\to0\) 与 \(t\to\infty\) 的两端都控制；将 \(r\geq1\) 代入后，剩余积分有统一上界或至多一个固定多项式因子，这个因子可再吸收到较慢的指数 \(e^{-cr}\) 中。因此取得远场指数衰减。
\end{proof}
这个核解释了非齐次空间的低频项：常数频率附近的符号有界，所以 \(J^{-s}\) 保持同一 \(L^p\)；齐次符号 \(|\xi|^{-s}\) 在零频发散，改善可积性时必须改变目标指数或规定低频归一化。两者在高频上具有相同阶数，却不能在全空间上互相替代。

\ProseParagraphBreak

核估计还给出 Sobolev 嵌入的另一种表述。若 \(1<p<\infty\)、\(0<s<d/p\)，置 \(1/q=1/p-s/d\)。由 \(G_s(x)\leq C|x|^{s-d}\) 和 HLS，
\[
 \|f\|_q=\|G_s*(J^sf)\|_q\leq C\|J^sf\|_p.
\]
若 \(s>d/p\)，则 \(G_s\in L^{p'}\)：原点幂次条件正是 \((d-s)p'<d\)，临界对数在有限指数也可积，远场由指数衰减处理。因此
\(\|f\|_\infty\leq\|G_s\|_{p'}\|J^sf\|_p\)，并得到连续代表。连续性由 \(L^{p'}\) 中的核平移连续性给出，而且紧支集输入的逼近说明该代表在无穷远趋零。

\ProseParagraphBreak

在临界 \(s=d/p\) 时，核不属于 \(L^{p'}\)，不能原样保留无穷端点。取
\(h(x)=|x|^{-s}(\log(e/|x|))^{-\beta}\mathbf1_{\{|x|<1/2\}}\)，其中 \(1/p<\beta\leq1\)。径向积分给出 \(h\in L^p\)；正核下界则使 \(G_s*h(0)\) 的径向积分成为 \(\int_0^{1/2}r^{-1}(\log(e/r))^{-\beta}\,\mathrm dr=\infty\)。对 \(x\to0\)，Fatou 给出输出任意大的邻域值，故它没有本性有界的代表。这个反例不仅指出一个点的发散，也排除了通过修改零测集恢复 \(L^\infty\) 的可能。

\ProseParagraphBreak

最后，椭圆导数估计与这些嵌入处理不同问题。前者从方程右端恢复同指数的最高阶导数，后者把已经拥有的导数换成更高的可积指数或连续性。应用偏微分方程时，常见路线是先证明解处于 \(H^{s,p}\)，再按 \(s\) 与 \(d/p\) 的相对位置使用嵌入。若只看方程的阶数就声称解连续，会遗漏右端的可积指数及维数这两个决定性参数。

\ProseParagraphBreak

齐次导数与非齐次范数之间还有一个常用的精确比较。对 \(s>0\)、\(1<p<\infty\)，在 Schwartz 类上有
\[
 \|J^sf\|_p\asymp\|f\|_p+\|(-\Delta)^{s/2}f\|_p.
\]
正向用两个乘子：\(J^{-s}\) 以及
\((2\pi|\xi|)^s(1+4\pi^2|\xi|^2)^{-s/2}\)。
后者在原点不一定普通光滑，但定义域只要求 \(\xi\ne0\)，其尺度导数依然受控，故满足 Mihlin 条件。反向取低频截断 \(\chi\)，低频 \(J^s\chi\) 的核可积；高频符号
\[
 (1-\chi(\xi))\frac{(1+4\pi^2|\xi|^2)^{s/2}}{(2\pi|\xi|)^s}
\]
是零阶 Mihlin 符号，把它作用于齐次导数即可恢复高频 \(J^sf\)。两部分合并证明比较；连续延拓把 Schwartz 类的闭包识别为相同的图范数空间。

\ProseParagraphBreak

因此非齐次空间不是拒绝齐次导数，而是在其之外补上低阶 \(L^p\) 大小。高频上的两个定义等价，真正差别留在零频附近。使用能量方法时，可以在已说明的空间范围内自由改写这个图范数；使用全空间齐次逆、缩放不变估计或商掉多项式时，却必须保留齐次定义自己的归一化。一个范数比较不能自动解除所有齐次分布的低频歧义。

\subsection{椭圆型估计}
\label{b2:subsec:220303}

设 \(A=(a_{jk})\) 为实对称正定常矩阵，且
\(\nu|\xi|^2\le\xi^{\mathsf T}A\xi\le\Lambda|\xi|^2\)。
对于紧支撑光滑 \(u\)，令
\(Lu=-\sum a_{jk}\partial_j\partial_k u\)。则
\[
 \widehat{\partial_j\partial_k u}
 =-\frac{\xi_j\xi_k}{\xi^{\mathsf T}A\xi}\widehat{Lu},
 \qquad
 \|D^2u\|_p\le C_{d,p,\nu,\Lambda}\|Lu\|_p.
\]
商是零阶齐次符号，其球面导数由椭圆下界控制；Mihlin 定理给出结论。这是先验估计，既不要求 \(L^{-1}f\) 在 \(L^p\) 中存在，也不将低频问题忽略。

\ProseParagraphBreak
若改用 \(1+L\)，符号
\((1+4\pi^2\xi^{\mathsf T}A\xi)^{-1}\) 及其乘上 \(\xi_j\xi_k\) 的版本均满足零阶界，所以
\[
 \|u\|_p+\|D^2u\|_p\le C\|(1+L)u\|_p.
\]
这可以扩展为 \((1+L)^{-1}:L^p\to W^{2,p}\) 的有界解算子，分布方程成立，唯一性在这一空间中可通过温和分布 Fourier 变换验证。

\ProseParagraphBreak
对于变系数方程，不能把 \(a_{jk}(x)\) 当常数移到 Fourier 符号中，因为 Fourier 变换不与空间乘法交换。局部冻结系数可以形成一个近似常系数问题，但需要小振荡条件、交换子估计和覆盖拼接。\chapref{b2:ch:16}的椭圆正则理论与\appref{b2:app:H}的低正则理论分别处理这些额外困难，本章的常系数证明不替代它们。

常系数模型可以进一步解释定义域与低频唯一性的角色。对正定矩阵 \(A\)，符号 \(1+4\pi^2\xi^{\mathsf T}A\xi\) 没有零点，它的逆及各阶导数至多多项式增长或衰减，因而乘法保持 Schwartz 测试空间。若温和分布 \(u\) 满足 \((1+L)u=0\)，把 Fourier 方程乘以该逆立即得到 \(\widehat u=0\)，故唯一性甚至在整个温和分布类中成立。这里有一个真正可逆的低频项。

\ProseParagraphBreak

对没有常数项的 \(L\)，情况不同。任何常数乃至适当的调和多项式都是齐次方程的温和分布解；若以 \(D^2u\) 的先验界工作，这些低阶自由度不会被看到。在 \(L^p(\mathbb R^d)\)、有限 \(p\) 中，非零常数当然不属于空间，但是否存在具有所需衰减的逆解，仍要检查频率原点的奇性。二阶导数估计已经给出，并不意味着 \(L^{-1}\) 自动为同指数有界算子。

\ProseParagraphBreak

在对角例子 \(L=-\sum_j a_j\partial_j^2\)、\(a_j>0\) 中，换元 \(x_j=\sqrt{a_j}\,y_j\) 将它化为标准 \(-\Delta_y\)。各个导数却分别得到 \(a_j^{-1/2}\) 或 \(a_j^{-1}\) 因子，因此当某个 \(a_j\downarrow0\) 时，椭圆正则估计会退化。一般正定矩阵可先正交对角化，正交变换保持 Lebesgue 测度与欧氏导数长度；最小特征值控制的正是这种退化。记录 \(\nu,\Lambda\) 的依赖有实际意义，而不是可删去的形式装饰。

\ProseParagraphBreak

带参数的预解算子把同一估计组织成一个尺度一致的族。对 \(\lambda>0\)，令 \(u=(\lambda+L)^{-1}f\)。以 \(\xi=\sqrt\lambda\,\zeta\) 归一化后，三个符号
\[
 \frac{\lambda}{\lambda+4\pi^2\xi^{\mathsf T}A\xi},\quad
 \frac{\sqrt\lambda\,2\pi\mathrm i\xi_j}{\lambda+4\pi^2\xi^{\mathsf T}A\xi},\quad
 \frac{-4\pi^2\xi_j\xi_k}{\lambda+4\pi^2\xi^{\mathsf T}A\xi}
\]
都具有与 \(\lambda\) 无关的 Mihlin 常数。于是
\[
 \lambda\|u\|_p+\sqrt\lambda\|\nabla u\|_p+\|D^2u\|_p
 \leq C_{d,p,\nu,\Lambda}\|f\|_p.
\]
符号缩放保持同指数算子范数，这是统一性的关键；直接逐项对未归一化符号求导往往会出现不必要的 \(\lambda^{-1}\) 因子。这里导数和零阶项分别以自己的尺度加权，也正好满足微分阶数的平衡。

\ProseParagraphBreak

令 \(\lambda\downarrow0\) 时，最高阶导数界仍统一，而对 \(\|u\|_p\) 的估计允许增长为 \(\lambda^{-1}\)。因此它没有绕过齐次逆算子的低频障碍。若需要在零参数得到实际解，必须另外控制数据的低频或选择齐次空间；有界参数预解族并不自动给出一个同指数的无质量逆。

\subsection{球面乘子 预告}
\label{b2:subsec:220304}

径向符号只依赖 \(|\xi|\)，仍可能在球面上发生跳跃。
光滑环带截断 \(\chi(|\xi|)\) 的核属于 \(L^1\)；球特征函数 \(\mathbf1_{\{|\xi|\le1\}}\) 则不满足任何跨边界的 Mihlin 光滑条件。
把跳跃平滑成宽度 \(\epsilon\) 的过渡层时，第 \(N\) 阶导数通常增长如 \(\epsilon^{-N}\)。对每个固定 \(\epsilon\) 应用定理可得有界性，却得不到 \(\epsilon\downarrow0\) 的统一界。

\ProseParagraphBreak
Bochner–Riesz 平均以 \((1-|\xi|^2)_+^\delta\) 逐渐消减边界跳跃，其所需光滑阶数与 \(p\)、维数有关。Mihlin 定理给出方便的充分条件，但常不是锐利范围。\chapref{b2:ch:24}用径向核的振荡和曲率研究这一问题；\chapref{b2:ch:23}的限制估计控制频率薄壳上的集中，则说明为什么乘子问题会与曲面延拓相遇。

\ProseParagraphBreak
因此“径向”“有界”“紧支撑”三项都不能单独替代乘子条件。对核使用 Young 不等式、对符号使用尺度导数、对薄壳使用曲面几何，是三条不同证明路线。应用时应明确正在使用哪一条，以及它在端点或退化极限中会失去哪一个条件。

平滑乘子还有一个时间演化模型，可与频率硬截断比较。热算子的符号是 \(e^{-4\pi^2t|\xi|^2}\)。其空间核非负，积分为一，因而在全部 \(L^p\) 上压缩；对每个整数 \(k\geq0\)，对 Gaussian 核求导及换元给出
\[
 \|D^ke^{t\Delta}f\|_p\leq C_k t^{-k/2}\|f\|_p,
 \qquad 1\leq p\leq\infty,\quad t>0.
\]
导数核是一个多项式乘 Gaussian，其 \(L^1\) 范数有限；时间缩放产生 \(t^{-k/2}\)，故该证明连两个端点也包括。单纯的开区间 Mihlin 定理已经足以处理 \(1<p<\infty\)，但核方法在这个特定模型中提供了额外端点。

\ProseParagraphBreak

任意实 \(s\geq0\) 时同样有
\[
 \|e^{t\Delta}f\|_{H^{s,p}}\leq C_{s,p}(1+t^{-s/2})\|f\|_p,
 \qquad1<p<\infty.
\]
在 \(0<t\leq1\) 中，把 \(t^{s/2}(1+4\pi^2|\xi|^2)^{s/2}e^{-4\pi^2t|\xi|^2}\) 以 \(\zeta=\sqrt t\,\xi\) 重写，成为
\((t+4\pi^2|\zeta|^2)^{s/2}e^{-4\pi^2|\zeta|^2}\)。
各个尺度导数在原点附近由 \(|\zeta|^{|\alpha|}\) 抵消可能的负幂，在无穷处由 Gaussian 控制，所以 Mihlin 常数统一；\(t\geq1\) 则用半群性将一次固定时间的平滑与剩余压缩复合。这个推导说明分数阶平滑并不要求先把 \(s\) 取整。

\ProseParagraphBreak

当 \(t\downarrow0\)，正阶范数常数发散，与一般 \(L^p\) 数据可能没有任何正阶导数完全相容。对每个固定正时间得到光滑函数，不意味着这些函数在更强空间中一致趋向原数据。连续半群的初值收敛发生在原来给定的空间，提升正则性则付出时间尺度。

\ProseParagraphBreak

把时间倒过来，符号成为 \(e^{4\pi^2t|\xi|^2}\)，甚至失去平方乘子的本性有界性。取 \(\widehat g\) 支于单位球，并令 \(f_N(x)=e^{2\pi\mathrm iNx_1}g(x)\)。输入平方范数不变，而逆热像的平方范数至少为
\(e^{4\pi^2t(N-1)^2}\|g\|_2\)。适当缩小输入振幅，可以使初值在 \(L^2\) 中趋零而逆演化的像趋于无穷。这个明确的高频检验说明正向平滑不能通过形式上的指数取反得到稳定逆过程。

\ProseParagraphBreak

与球面乘子相比，热符号没有一条频率边界上的跳跃，所有尺度的导数也具有统一控制；与齐次分数积分相比，它在零频不发散。良好的端点压缩、正时间平滑和低频稳定是这些特殊性质共同产生的，不应当归给任意一个写成 \(m(\xi)\) 的符号。

\ProseParagraphBreak

频率过渡层还可以用一个明确的算子范数检验。取固定 \(\chi\in C_c^\infty\)，令 \(m_R(\xi)=\chi(\xi/R)\)。其核为 \(R^d\check\chi(Rx)\)，所以全部 \(R>0\) 的 \(L^1\) 核范数都等于 \(\|\check\chi\|_1\)，在每个 \(L^p\) 上有同一个界。若 \(\chi(0)=1\)，在 \(R\to\infty\) 时，有限指数下这族算子强趋向恒等：先对 Schwartz 输入用核的近似恒等性，再用一致范数和稠密性推广。将频率支撑推向无穷本身并不造成常数发散；关键是过渡轮廓经过尺度归一化后仍固定。

\ProseParagraphBreak

若改用 \(m_\epsilon(\xi)=\theta((|\xi|-1)/\epsilon)\)，其中 \(\theta\) 固定而过渡层紧贴单位球面，归一化整个频率空间不能同时固定球的半径和层的厚度。径向一阶导数在该层上达到 \(\epsilon^{-1}\)，而 \(|\xi|\) 仍约为一，尺度导数条件因此真实地增长。两个模型都叫作光滑截断，却有不同的无量纲参数：前者只有整体半径，后者还有厚度与半径之比。将其混同，会把 Mihlin 常数的一致性误判成一个无条件事实。

\ProseParagraphBreak

在平方空间中只需符号的本性上确界，故上述有界过渡层即便变得越来越薄，仍可具有统一平方界。对其他指数则必须利用核、导数或特殊几何。这个差别再次说明，平方乘子的全部刻画并没有自动延伸到任意 \(L^p\)。讨论逼近一个不连续乘子时，要分别检查固定参数有界、一致参数有界和极限算子的有界性；只有前两项共同存在，才可能通过稠密输入上的极限传递到后一项。
